\chapter{Previsión de la demanda}
\section{Papel de las previsiones. Principios de previsión de la demanda}
Se conoce como previsión de demanda a la proyección de la demanda esperada en el futuro, dadas unas determinadas condiciones de entorno. La previsión de la demanda es el punto de partida de la planificación en las redes de suministro.

Principios básicos en previsión de la demanda:
\begin{itemize}
    \item Las previsiones siemmpre tienen un error asociado. El error dependerá del modelo y de la naturaleza de los datos.
    \item Cuanto mayor es la agregación, mejor es la previsión.
    \item A mayor horizonte, peor es la previsión (para mismo nivel de agregación). Las previsiones a largo plazo de la demanda diaria presentan un error relativo mayor que las previsiones a corto plazo de esa misma demanda diatia (menor fiabilidad).
    \item El horizonte temporal y el nivel de agregación deben elegirse de forma consistente. Las previsiones a largo plazo requieren un mayor nivel de agregación. Las previsiones a corto plazo pueden permitir un menor nivel de agregación.
\end{itemize}

\section{Factores a considerar en la selección de técnicas de previsión}


\section{Tipos de demanda}

\section{Métodos de previsión}
\subsection{Enfoques cualitativos}
Se basan en opiniones, de ahí el carácter subjetivo de la información de partida, aunque también pueden incorporar datos históricos. Recogen de forma sistemática toda la información. Normalmente, la previsión resultante se cuantifica usando un cierto tratamiento matemático.
Se usan cuando:
\begin{itemize}
    \item Previsiones a largo plazo o de largo alcance
    \item No hay datos históricos o no son fiables
\end{itemize}

Ejemplo:
\begin{itemize}
    \item Estudios de mercado
    \item Método Delphi (consulta iterativa a expertos)
    \item Analogía del ciclo de vida
    \item Agregación de opiniones / experiencia de la fuerza de ventas
\end{itemize}

\subsection{Enfoques cuantitativos}
Basados en datos históricos y/o variables causales. 
Análisis de series temporales:
\begin{itemize}
    \item Usan datos históricos de ventas para hacer la previsión
    \item Se basan en la hipótesis de que la demanda del pasado es un buen indicador de la demanda futura
\end{itemize}

Modelos causales:
\begin{equation}
    P(t) = f(X_t)
\end{equation}

\section{Modelos de series temporales}
Concepto: 
\begin{itemize}
    \item Conjunto de observaciones de los valores de una variable a lo largo del tiempo.
\end{itemize}

Hipótesis:
\begin{itemize}
    \item La demanda prevista en el futuro es una función de la demanda observada en el pasado.
    \item La demanda se puede dividir en componentes.
\end{itemize}

Componentes:
\begin{itemize}
    \item Nivel de referencia o estacionario
    \subitem Escala de la serie. Si sólo existe este componente, la serie permanece constante.
    \item Tendencia
    \subitem Tasa de crecimiento o declive de la serie a lo largo del tiempo (no necesariamente lineal).
    \item Estacionalidad
    \subitem Patrón cíclico de incrementos y decrecimientos de la demanda.
    \subitem Generado por clima, decisiones de clientes.
    \subitem Pueden superponerse distintas estacionalidades
    \item Componente aleatorio
    \subitem Fluctuación aleatoria que no puede ser explicada por la historia de ventas.
\end{itemize}

\subsection{Media móvil}

\subsection{Alisado exponencial simple}

\subsection{Alisado exponencial doble (Holt)}

\subsection{Método de Winters}

\section{Medidas del error de previsión y señal de seguimiento}